\section{Concluding remarks}
\label{sec:conclusion}

The signed spectral problem on $C_n(1,2)$ is governed more naturally by cycle
flux than by individual edge signs.  Negative half-cell flux yields a general
monomial chiral involution and halves the squared Bloch problem.  Period eight
is the first primitive period at which this mechanism produces squared edge
below $8$, and the resulting antipodal two-defect orbit is unique at that
period.  Its additional block structure makes the entire dispersion explicit,
so both finite holonomy sectors have exact spectral radii.  The positive sector
then gives an infinite family that strictly improves the twisted benchmark.

For finite rings, the negative-holonomy formula
\eqref{eq:negative-finite-exact} gives the stronger bound
\[
 m(C_{8L}(1,2))\leq\rho(A_{8L,-}^{*})<\sqrt\etaedge
 \qquad(L\geq1).
\]
The remaining extremal problem is to determine $m(C_{8L}(1,2))$ and its
minimizing switching classes.  The negative-holonomy phase provides an explicit
upper bound against which other periodic and nonperiodic signings can be
compared.  The half-cell criterion and the defect obstructions supply
structural tools for this comparison.
